\usepackage[T1]{fontenc}
\usepackage{lmodern}
\pdfmapfile{+lm.map}
\usepackage[utf8]{inputenc}
\usepackage{amsmath,amssymb,amsthm,mathtools}
\usepackage{mathrsfs}
\usepackage{microtype}
\usepackage[margin=1.03in]{geometry}
\usepackage{graphicx,booktabs,array}
\usepackage{xcolor}
\definecolor{linkblue}{RGB}{28,68,110}
\usepackage[colorlinks=true,linkcolor=linkblue,citecolor=linkblue,urlcolor=linkblue,bookmarksnumbered=true]{hyperref}
\hypersetup{pdftitle={An L3 bound for the trilinear Hilbert transform},pdfauthor={OpenAI}}
\newtheorem{theorem}{Theorem}[section]
\newtheorem{lemma}[theorem]{Lemma}
\newtheorem{proposition}[theorem]{Proposition}
\newtheorem{corollary}[theorem]{Corollary}
\theoremstyle{definition}
\newtheorem{definition}[theorem]{Definition}
\theoremstyle{remark}
\newtheorem{remark}[theorem]{Remark}
\newcommand{\R}{\mathbb R}
\newcommand{\C}{\mathbb C}
\newcommand{\Z}{\mathbb Z}
\newcommand{\Q}{\mathbb Q}
\newcommand{\N}{\mathbb N}
\newcommand{\T}{\mathbb T}
\newcommand{\E}{\mathbb E}
\newcommand{\Prob}{\mathbb P}
\newcommand{\1}{\mathbf 1}
\newcommand{\ip}[2]{\langle #1,#2\rangle}
\newcommand{\norm}[1]{\lVert #1\rVert}
\newcommand{\abs}[1]{\lvert #1\rvert}
\newcommand{\phys}{\mathrm{phys}}
\newcommand{\fint}{\mathop{\mathchoice{\mkern2mu\int\mkern-15mu-\mkern7mu}{\mkern2mu\int\mkern-13mu-\mkern6mu}{\mkern2mu\int\mkern-11mu-\mkern5mu}{\mkern2mu\int\mkern-9mu-\mkern4mu}}\nolimits}
\DeclareMathOperator{\supp}{supp}
\DeclareMathOperator{\dist}{dist}
\DeclareMathOperator{\diam}{diam}
\DeclareMathOperator{\rank}{rank}
\DeclareMathOperator{\spanop}{span}
\DeclareMathOperator{\sgn}{sgn}
\setlength{\emergencystretch}{1.5em}
\allowdisplaybreaks[2]
\pdfgentounicode=1
\input{glyphtounicode}
